\subsection{Fixed-response activation and explicit cusp calculations}

\paragraph{Onset with fixed responses.}
The next result varies the coefficient of the polynomials $P_i$ rather
than the shape of the responses.

\beT[Increasing higher-order terms creates positive equilibria
while zero remains stable][activation by higher-order terms][ActF]\label{t:KBactivation}
Consider
\[
 \dot x=f(x;s)=-Vx+\Pi_s(x),\qquad
 (\Pi_s)_i=b_i\Phi_i(sP_i(x)),\qquad s\ge0,
\]
where $V$ is a nonsingular $M$-matrix, $b\gg0$, and the invariant
state space is $[0,\bar x]$, $\bar x=V^{-1}b$.
Each $P_i$ is a nonzero polynomial with nonnegative coefficients and
monomials of degree at least two. Assume
$\Phi_i\in C^1([0,\infty))$, with the derivative at zero understood from
the right, and
\[
 \Phi_i(0)=0,\qquad \Phi_i(t)\uparrow1\quad(t\to\infty),
 \qquad \Phi_i'(t)>0\quad(t>0).
\]
The graph containing transfer edges $j\to i$ with $-v_{ij}>0$ and
edges $j\to i$ supplied by pure-power terms $a x_j^k$ in $P_i$, $a>0$,
is assumed strongly connected. Mixed monomials are allowed without
supplying these edges. Then $s_*\in(0,\infty)$, defined by
\begin{equation}\label{eq:KB-activation}
 s_*:=\inf\{s\ge0:\exists x\gg0,\ x\le\bar x,\ \Pi_s(x)\ge Vx\},
\end{equation}
is attained by a positive equilibrium. Positive equilibria exist
exactly for $s\ge s_*$, while zero is locally exponentially stable
for every finite $s$.

At the greatest equilibrium $x_*$ for $s=s_*$, the Jacobian has a
simple zero eigenvalue and a strictly stable complement. If $f$ is
$C^2$ near $(x_*,s_*)$ and
\[
 \ell^\top\partial_x^2f(x_*;s_*)[r,r]\ne0,
 \qquad Jr=0,\quad\ell^\top J=0,\quad r,\ell\gg0,
 \quad\ell^\top r=1,
\]
then, for $s>s_*$ sufficiently close to $s_*$, a saddle-node creates
a locally exponentially stable positive equilibrium and a positive
equilibrium with exactly one unstable eigenvalue.
\eeT
\begin{proof}
A positive subsolution $\Pi_s(x)\ge Vx$ generates a nondecreasing
trajectory bounded by $\bar x$, whose limit is a positive equilibrium.
Increasing $s$ preserves the subsolution inequality. For $0<t<1$,
$V(t\bar x)=tb$ and every $P_i(t\bar x)>0$, so $t\bar x$ is a
subsolution for sufficiently large $s$. For sufficiently small $s$,
\[
 \sup_{x\in[0,\bar x]}\|V^{-1}\partial_x\Pi_s(x)\|_\infty<1,
\]
since $\partial_x\Pi_s\to0$ uniformly on the box. Thus
$x\mapsto V^{-1}\Pi_s(x)$ is a contraction of the box into itself;
its fixed point is zero. Consequently $0<s_*<\infty$.
On bounded $s$-intervals, $\Pi_s(x)=O(\|x\|^2)$ uniformly near zero;
$x=V^{-1}\Pi_s(x)$ therefore excludes small nonzero equilibria.
Limits as $s\downarrow s_*$ give a nonzero equilibrium at $s_*$.
At an equilibrium, a zero coordinate forces its graph predecessors to
be zero; strong connectivity makes every nonzero equilibrium positive.
Also $J(0)=-V$, giving stability of zero.

The decreasing trajectory from $\bar x$ gives the greatest equilibrium
$x_*$. It lies strictly below $\bar x$, since
\[
 \bar x-x_*=V^{-1}(b-\Pi_{s_*}(x_*))\gg0.
\]
Its Jacobian is irreducible Metzler. If $s(J)>0$, the positive Perron
vector gives a strict subsolution $x_*+\varepsilon r$ inside the same
box, contradicting maximality. If $s(J)<0$, the implicit-function
theorem continues the positive equilibrium below $s_*$. Hence $s(J)=0$
is simple and all other eigenvalues have negative real parts.
Finally,
\[
 \ell^\top\partial_s f(x_*;s_*)
 =\sum_i\ell_i b_i\Phi_i'(s_*P_i(x_*))P_i(x_*)>0.
\]
The saddle-node test of Proposition~\ref{p:sadCus} applies when
its quadratic coefficient is nonzero.
That coefficient is negative because positive equilibria do not exist
below $s_*$. The two branches above it have zero and one unstable
eigenvalues, respectively; the latter is a saddle in dimension at least two.
\end{proof}

\paragraph{Homogeneous branches, saddle-nodes, and heterogeneous cusps.}

\beT[Balanced response arguments preserve scalar branches and their first instability]
[homogeneous branches][KBhom]
For \eqref{eq:KB}, suppose
$V\mathbf1=v\mathbf1$, $v>0$, and
$\Pi(s\mathbf1)=r(s)\mathbf1$ throughout the relevant interval.
The states $x=s\mathbf1$, in which all variables are equal, then obey
\begin{equation}\label{eq:KBscalar}
 \dot s=-vs+r(s).
\end{equation}
If the Jacobian $J$ at such an equilibrium is Metzler, its spectral
abscissa is $s(J)=-v+r'(s)$.
When this number is zero the eigenvalue zero is semisimple, even
without irreducibility. If $J$ is irreducible it is simple, and all
heterogeneous eigenvalues have negative real parts.
For $r(s)=b\Phi(h+\beta s)$ with $b,\beta>0$,
write $\varphi=\Phi'$ and assume the derivatives used below exist and
are continuous. A scalar saddle-node candidate
satisfies
\[
 vs=b\Phi(z),\quad z=h+\beta s,\quad b\beta\varphi(z)=v.
\]
It is nondegenerate for variation of $h$ if $\varphi'(z)\ne0$.
If $\varphi(z_c)>0$, $\varphi'(z_c)=0$, and
$\varphi''(z_c)\ne0$, the scalar equation has a nondegenerate two-parameter cusp at
\[
 s_c=\frac b v\Phi(z_c),\qquad
 \beta_c=\frac v{b\varphi(z_c)},\qquad
 h_c=z_c-\beta_cs_c.
\]
At a saddle-node or cusp with irreducible Metzler full Jacobian, the scalar
bifurcation has only stable complementary directions in the full system.
For higher-order balanced $P_i$, one may take
$r(s)=b\Phi(h+\sum_k a_ks^k)$ instead.
\eeT
\begin{proof}
Invariance follows by substitution. With $H=\mathbf1$,
Theorem~\ref{t:Spe}(1) gives spectral equality. More explicitly,
$J\mathbf1=(-v+r'(s))\mathbf1$ implies
$\|e^{Jt}\|_\infty=e^{(-v+r'(s))t}$. This proves spectral ordering
and, at zero, boundedness of the semigroup excludes Jordan blocks.
The scalar derivatives are
$f_s=-v+b\beta\varphi$, $f_{ss}=b\beta^2\varphi'$, and
$f_{sss}=b\beta^3\varphi''$. At the cusp the unfolding determinant
in $(h,\beta)$ is $b^2\varphi(z_c)^2>0$.
Proposition~\ref{p:sadCus} verifies these scalar bifurcations.
The cusp or saddle-node lifts under Theorem~\ref{t:Spe}(3); in a
reducible network several center directions may instead coexist.
\end{proof}

For $V=I$ and row-stochastic $W$, the spectral ordering and the semisimple or simple zero
eigenvalue are in \cite{Boud26}; Theorem~\ref{t:KBhom} extends that argument to
$V\mathbf1=v\mathbf1$ and nonlinear balanced production.

For a unimodal density with $\varphi''(z_c)<0$, the cubic coefficient $f_{sss}$ is negative, so the outer branches of the cusp
attract, and for $\beta>\beta_c$ the two folds bound a hysteresis loop in $h$ (Figure~\ref{fig:hyst}).

\begin{figure}[H]
\centering
\begin{minipage}{.48\linewidth}
\centering
\begin{tikzpicture}
\begin{axis}[width=\linewidth,height=.78\linewidth,xlabel={$h$},ylabel={$s$},ymin=0,ymax=1,
 tick label style={font=\scriptsize},label style={font=\small},title style={font=\small},xmin=-1.9,xmax=0.4,title={$\beta=1.5<\beta_c$}]
\addplot[blue,thick] coordinates {(-3.2010,0.0007) (-3.0418,0.0012) (-2.8830,0.0020) (-2.7249,0.0033) (-2.5679,0.0052) (-2.4123,0.0082) (-2.2588,0.0125) (-2.1081,0.0188) (-1.9611,0.0274) (-1.8188,0.0392) (-1.6822,0.0548) (-1.5524,0.0749) (-1.4304,0.1003) (-1.3170,0.1314) (-1.2128,0.1685) (-1.1178,0.2119) (-1.0316,0.2611) (-0.9534,0.3156) (-0.8817,0.3745) (-0.8147,0.4364) (-0.7500,0.5000) (-0.6853,0.5636) (-0.6183,0.6255) (-0.5466,0.6844) (-0.4684,0.7389) (-0.3822,0.7881) (-0.2872,0.8315) (-0.1830,0.8686) (-0.0696,0.8997) (0.0524,0.9251) (0.1822,0.9452) (0.3188,0.9608) (0.4611,0.9726) (0.6081,0.9812) (0.7588,0.9875) (0.9123,0.9918) (1.0679,0.9948) (1.2249,0.9967) (1.3830,0.9980) (1.5418,0.9988) (1.7010,0.9993)};
\end{axis}
\end{tikzpicture}
\end{minipage}\hfill
\begin{minipage}{.48\linewidth}
\centering
\begin{tikzpicture}
\begin{axis}[width=\linewidth,height=.78\linewidth,xlabel={$h$},ylabel={$s$},ymin=0,ymax=1,
 tick label style={font=\scriptsize},label style={font=\small},title style={font=\small},xmin=-3.2,xmax=-0.8,title={$\beta=4>\beta_c$}]
\addplot[blue,thick] coordinates {(-3.8003,0.0001) (-3.7296,0.0001) (-3.6588,0.0001) (-3.5882,0.0002) (-3.5176,0.0002) (-3.4470,0.0003) (-3.3765,0.0004) (-3.3061,0.0005) (-3.2358,0.0006) (-3.1657,0.0008) (-3.0957,0.0010) (-3.0259,0.0013) (-2.9564,0.0016) (-2.8872,0.0020) (-2.8183,0.0025) (-2.7499,0.0031) (-2.6820,0.0038) (-2.6148,0.0047) (-2.5482,0.0058) (-2.4825,0.0071) (-2.4177,0.0086) (-2.3541,0.0104) (-2.2917,0.0125) (-2.2308,0.0150) (-2.1715,0.0179) (-2.1141,0.0212) (-2.0588,0.0251) (-2.0058,0.0295) (-1.9553,0.0346) (-1.9076,0.0404) (-1.8629,0.0470) (-1.8216,0.0543) (-1.7838,0.0626) (-1.7498,0.0718) (-1.7198,0.0820) (-1.6940,0.0933) (-1.6726,0.1056) (-1.6559,0.1191) (-1.6438,0.1338) (-1.6365,0.1497) (-1.6341,0.1668)};
\addplot[blue,thick] coordinates {(-2.3659,0.8332) (-2.3635,0.8503) (-2.3562,0.8662) (-2.3441,0.8809) (-2.3274,0.8944) (-2.3060,0.9067) (-2.2802,0.9180) (-2.2502,0.9282) (-2.2162,0.9374) (-2.1784,0.9457) (-2.1371,0.9530) (-2.0924,0.9596) (-2.0447,0.9654) (-1.9942,0.9705) (-1.9412,0.9749) (-1.8859,0.9788) (-1.8285,0.9821) (-1.7692,0.9850) (-1.7083,0.9875) (-1.6459,0.9896) (-1.5823,0.9914) (-1.5175,0.9929) (-1.4518,0.9942) (-1.3852,0.9953) (-1.3180,0.9962) (-1.2501,0.9969) (-1.1817,0.9975) (-1.1128,0.9980) (-1.0436,0.9984) (-0.9741,0.9987) (-0.9043,0.9990) (-0.8343,0.9992) (-0.7642,0.9994) (-0.6939,0.9995) (-0.6235,0.9996) (-0.5530,0.9997) (-0.4824,0.9998) (-0.4118,0.9998) (-0.3412,0.9999) (-0.2704,0.9999) (-0.1997,0.9999)};
\addplot[blue,dashed] coordinates {(-1.6341,0.1668) (-1.6372,0.1877) (-1.6465,0.2102) (-1.6619,0.2342) (-1.6830,0.2596) (-1.7095,0.2864) (-1.7410,0.3144) (-1.7770,0.3435) (-1.8168,0.3736) (-1.8597,0.4045) (-1.9051,0.4360) (-1.9521,0.4679) (-2.0000,0.5000) (-2.0479,0.5321) (-2.0949,0.5640) (-2.1403,0.5955) (-2.1832,0.6264) (-2.2230,0.6565) (-2.2590,0.6856) (-2.2905,0.7136) (-2.3170,0.7404) (-2.3381,0.7658) (-2.3535,0.7898) (-2.3628,0.8123) (-2.3659,0.8332)};
\draw[red,thick,-{Stealth[length=5pt]}] (axis cs:-1.6341,0.1668) -- (axis cs:-1.6341,0.9900);
\draw[red,thick,-{Stealth[length=5pt]}] (axis cs:-2.3659,0.8332) -- (axis cs:-2.3659,0.0100);
\end{axis}
\end{tikzpicture}
\end{minipage}
\caption{Hysteresis on the homogeneous branch of Theorem~\ref{t:KBhom}, $\dot s=-s+\Phi_{\mathrm G}(h+\beta s)$, for the Gaussian response:
the equilibria are $s=\Phi_{\mathrm G}(z)$, $h=z-\beta\Phi_{\mathrm G}(z)$, and $\beta_c=\sqrt{2\pi}$. Left: the response to $h$ is
gradual and reversible. Right: solid branches are stable and the dashed branch is unstable. As $h$ increases the state jumps up at
$h\approx-1.634$; as $h$ decreases it jumps down at $h\approx-2.366$.}
\label{fig:hyst}
\end{figure}

\beT[Two identical response populations admit cusp formulas with linear exchange]
[cusps with linear exchange][KBexchange]
Consider $\dot x=-V_\kappa x+\Pi(x)$ for two populations with
the same strictly increasing response $\Phi:\mathbb R\to(0,1)$ and coefficients,
\[
 V_\kappa=\begin{pmatrix}v+\kappa&-\kappa\\-\kappa&v+\kappa\end{pmatrix},
 \quad W_\varepsilon=\begin{pmatrix}1-\varepsilon&\varepsilon\\
 \varepsilon&1-\varepsilon\end{pmatrix},
 \quad \Pi_i=b\Phi(h+\beta(W_\varepsilon x)_i),
\]
where $v,b,\beta>0$, $\kappa\ge0$, and
$0\le\varepsilon\le1/2$. The terms with coefficient $\kappa$
transfer activity between the two variables and cancel in their sum.
Write the two response arguments as $a+s$ and $a-s$, and put
\[
 M(a,s)=\tfrac12(\Phi(a+s)+\Phi(a-s)),\quad
 D(a,s)=\tfrac12(\Phi(a+s)-\Phi(a-s)),\quad
 \widehat\beta=\beta b/v.
\]
The heterogeneous equilibria are parameterized, whenever $D\ne0$, by
\begin{align*}
 m&=\frac b v M,&d&=\frac b{v+2\kappa}D,&x_{1,2}&=m\pm d,\\
 h&=a-\widehat\beta M,&
 \varepsilon&=\frac12\left(1-
 \frac{(v+2\kappa)s}{\beta bD}\right).
\end{align*}
At the same arguments $(a,s)$, define the parameter map without exchange,
with unit removal and production coefficient, by
\[
 U=a-\widehat\beta M,\qquad
 \varepsilon_0=\frac12\left(1-\frac{s}{\widehat\beta D}\right).
\]
The original parameters satisfy exactly
\begin{equation}\label{eq:KB-cusp-affine}
 (h,\varepsilon)=
 \left(U,(1+2\kappa/v)\varepsilon_0-\kappa/v\right).
\end{equation}
Consequently its generic parameter-map folds and cusps are preserved
by this affine transformation, retaining only those with admissible
$0\le\varepsilon\le1/2$ and positive activities.
Suppose $\Phi$ is $C^7$ near $z_c$ and, writing $\varphi=\Phi'$,
\[
 \varphi(z_c)>0,\qquad \varphi'(z_c)=0,\qquad
 \varphi''(z_c)<0.
\]
Put $a_0=b/v$, $\beta_c=1/(a_0\varphi(z_c))$,
$x_c=a_0\Phi(z_c)$, $h_c=z_c-\beta_cx_c$, and
$\delta=\beta-\beta_c\downarrow0$.
The two cusps related by interchanging the populations extend the
zero-exchange calculation of \cite{Boud26}. They satisfy,
for $\kappa=O(\delta)$,
\begin{align*}
 \varepsilon_c&=\frac{\delta}{3\beta_c}-\frac\kappa v+o(\delta),\\
 |x_{1,c}-x_{2,c}|&=
 2\left[-\frac{2}{a_0\varphi''(z_c)\beta_c^4}\right]^{1/2}
 \delta^{1/2}+o(\delta^{1/2}),\\
 h_c(\beta)-h_c&=-x_c\delta+o(\delta).
\end{align*}
Let $\varepsilon_{0,c}(\widehat\beta)$ be the value of $\varepsilon_0$ at the cusps. Their
couplings are admissible exactly when $\varepsilon_{0,c}(\widehat\beta)\ge\kappa/(v+2\kappa)$.
For fixed $\kappa>0$, these cusps leave the physical range
$\varepsilon\ge0$ sufficiently near the scalar cusp.
\eeT
\begin{proof}
The mean equilibrium equation is $vm=bM$, and the contrast equation
is $(v+2\kappa)d=bD$. The fields satisfy
$a=h+\beta m$ and $s=\beta(1-2\varepsilon)d$.
Eliminating $m,d$ gives the parameterization and
\eqref{eq:KB-cusp-affine}. Its target transformation is invertible,
so it preserves Whitney fold and cusp type. The source transformation
from activities to fields is locally invertible when
$\beta(1-2\varepsilon)\ne0$.

\emph{Unit model.} Take $v=b=1$, $\kappa=0$, and write
$\delta_u=\beta-\beta_c$. This is the model of \cite{Boud26}, whose theorem on
heterogeneous cusps born from the scalar cusp gives two exchange-related generic Whitney
cusps of the parameter map $(a,s)\mapsto(h,\varepsilon)$ and the three displayed laws.
We record an exact form of its nondegeneracy computation.
Let $\Xi=\partial_ah\,\partial_s\varepsilon-\partial_sh\,\partial_a\varepsilon$
and let $D_k\Xi=\partial_a\Xi\,\partial_s\varepsilon-\partial_s\Xi\,\partial_a\varepsilon$
be its derivative along the kernel direction $k=(\partial_s\varepsilon,-\partial_a\varepsilon)$. Put
\[
 R(a,s)=\frac{D(a,s)}s
       =\frac12\int_{-1}^1\varphi(a+\tau s)\,d\tau,
\]
\[
 H=(1-\beta M_a)R_s+\beta M_sR_a,
 \qquad \mathcal K=R_s\partial_a-R_a\partial_s.
\]
Since $\mathcal K R=0$, direct differentiation gives the exact identities
\[
 \Xi=\frac{H}{2\beta R^2},\qquad
 D_k\Xi=\frac{\mathcal K H}{4\beta^2R^4},\qquad
 D_k^2\Xi=\frac{\mathcal K^2H}{8\beta^3R^6}.
\]
With $\delta_u=t^2$, $s=tS$, $a-z_c=t^2A$, the cusp branches have
\[
 S(0)=\pm S_*,\quad S_*^2=-\frac{2}{\beta_c^2\varphi''(z_c)},\qquad
 A(0)=\frac{\varphi'''(z_c)}{4\beta_c^2\varphi''(z_c)^2},
\]
and, with $\varphi''=\varphi''(z_c)$, along them
\[
 \partial_s\Xi=\frac{\varphi''}{3}t^2+o(t^2),\qquad
 \partial_s\varepsilon=\frac{\beta_c\varphi''S}{6}t+o(t),\qquad
 D_k^2\Xi=-\frac{2\beta_c^2\varphi''^3S}{27}t^3+o(t^3).
\]
Thus $d\Xi\ne0$, the parameter map has rank one, and
$D_k^2\Xi\ne0$ for all sufficiently small $t>0$.

\emph{General model.} The activities are $a_0=b/v$ times the unit-model
ones. With $\widehat\beta-\widehat\beta_c=a_0\delta$ and
$\widehat\beta_c=a_0\beta_c$, the unit-model formulas give the displayed ones
for $\kappa=0$; the transformation \eqref{eq:KB-cusp-affine} gives
$\varepsilon_c=(1+2\kappa/v)\varepsilon_{0,c}-\kappa/v$, and the contrast carries the
factor $b/(v+2\kappa)$ in place of $b/v$. For $\kappa=O(\delta)$, the product
$(2\kappa/v)\varepsilon_{0,c}$ is $O(\delta^2)$ and the correction
to the contrast is $O(\delta^{3/2})$; the leading coupling shift
is $-\kappa/v$.
By \eqref{eq:KB-cusp-affine}, $0\le\varepsilon\le1/2$ is equivalent to
$\kappa/(v+2\kappa)\le\varepsilon_0\le1/2$, which is the admissibility test of the statement.
Near the cusp the other two-dimensional eigenvalue is negative:
at the same fields, its trace satisfies exactly
\[
 \operatorname{tr}J_\kappa
 =v[-2+(1-\varepsilon_{0,c})(q_1+q_2)]
  +\kappa[(1-2\varepsilon_{0,c})(q_1+q_2)-2],
\]
where
$q_i=\widehat\beta\varphi(a\pm s)$.
Here $q_1+q_2=2+o(\delta)$ and
$\varepsilon_{0,c}=\delta/(3\beta_c)+o(\delta)$.
For $\kappa=O(\delta)$, the trace is consequently
$-2v\delta/(3\beta_c)+o(\delta)<0$.
Thus physical cusps have one stable complementary direction.
Theorem~\ref{t:Spe}(3) lifts them to an exact cooperative
network realization whose Jacobian at the cusp is irreducible.
\end{proof}

\paragraph{Application: persistence of the two exchange cusps.}
For each fixed sufficiently small $\beta-\beta_c>0$ for which
Theorem~\ref{t:KBexchange} supplies two distinct nondegenerate
full-network cusps, Proposition~\ref{p:sadCus} continues both under
sufficiently small $C^3$ asymmetry or loss of equitability.
Apply it in disjoint neighbourhoods of the two cusp states.
Their stimulus and coupling values need no longer coincide,
and the permissible perturbation size may depend on $\beta-\beta_c$.
Interior admissibility persists; at a parameter boundary such as
$\varepsilon=0$, check it separately using the exchange-induced
admissibility inequality above.

\beT[Unequal stimuli preserve both threshold cusps precisely within an explicit asymmetry interval]
[][cuspAsym]
Consider
\[
 \dot x_{1,2}=-x_{1,2}+
 \Phi\bigl(h\pm\eta+\beta[(1-k)x_{1,2}+kx_{2,1}]\bigr),
 \qquad 0\le k\le1/2,\quad\beta>0.
\]
Let $\Phi:\mathbb R\to(0,1)$ be a strictly increasing CDF, $C^5$ near zero,
with $\Phi(-q)=1-\Phi(q)$. Put $\varphi=\Phi'$, assume $\varphi(0)>0$, $\varphi''(0)<0$,
and set $\beta_c=1/\varphi(0)$. For $\beta=\beta_c+\delta$, $\delta>0$
sufficiently small, there is a unique sufficiently small $s>0$ with
$\beta\varphi(s)=1$. Define
\[
 D=\Phi(s)-\tfrac12>0,\qquad \eta_* =\beta D-s>0,\qquad
 k_0=\frac{\eta_*}{2\beta D}.
\]
For $\zeta=\pm1$, the two continuations of the heterogeneous cusps at
$\eta=0$ are
\[
 x_\zeta^*=(\tfrac12+\zeta D,\tfrac12-\zeta D),\qquad
 h_\zeta^*=-\beta/2,\qquad k_\zeta^*=k_0+\frac{\zeta\eta}{2\beta D}.
\]
Both candidates lie in $0<k<1/2$ and are nondegenerate cusps of
the $(h,k)$ family precisely while
$|\eta|<\eta_*(\beta)$, with eigenvalues $0,-2k_\zeta^*<0$.
At $|\eta|=\eta_*$ one reaches $k=0$ with two zero eigenvalues,
ceasing to be a one-dimensional-centre cusp; beyond it that candidate
has inadmissible $k<0$. The other cusp remains nondegenerate for
asymmetries near this boundary. Moreover,
\[
 s=\sqrt{\frac{-2}{\beta_c^2\varphi''(0)}}\delta^{1/2}+O(\delta^{3/2}),\qquad
 \eta_* =\frac{2}{3\beta_c}\sqrt{\frac{-2}{\beta_c^2\varphi''(0)}}
 \delta^{3/2}+O(\delta^{5/2}).
\]
\eeT
\begin{proof}
In coordinates $m=(x_1+x_2)/2-1/2$, $d=(x_1-x_2)/2$, at
$h=-\beta/2$, $m=0$ the arguments are $q_1=\eta+\beta(1-2k)d$, $q_2=-q_1$.
The mean equation vanishes by symmetry; setting $d=\zeta D$, $q_1=\zeta s$
gives $\zeta s=\eta+\beta(1-2k)\zeta D$, hence $k=k_\zeta^*$.
Evenness of $\varphi$ and $\beta\varphi(s)=1$ give
\[
 J=-I+\begin{pmatrix}1-k&k\\k&1-k\end{pmatrix},
\]
with eigenvalues $0,-2k$ and mean-direction kernel for $k>0$.

Write the equations as $f_m,f_d$. Since $\partial_df_d=-2k$, solve
$f_d=0$ for $d=d(m,h,k)$ and put $g=f_m(m,d(m,h,k),h,k)$.
At either candidate $d_m=0$. Differentiating $f_d=0$ twice gives
\[
 d_{mm}=\frac{\zeta\beta^2\varphi'(s)}{2k},\qquad
 d_k=-\frac{\zeta D}{k}.
\]
Substitution into the derivatives of $g$ gives
\[
 g=g_m=g_{mm}=0,\qquad
 g_{mmm}=\beta^3\varphi''(s)+\frac{3\beta^4(1-2k)\varphi'(s)^2}{2k},
\]
\[
 g_h=1/\beta,\quad g_k=0,\quad g_{mk}=-\frac{\beta^2D\varphi'(s)}k,
 \qquad
 \det\begin{pmatrix}g_h&g_k\\g_{mh}&g_{mk}\end{pmatrix}
 =-\frac{\beta D\varphi'(s)}k\ne0.
\]
These formulas apply to both signs because $\varphi'$ is odd;
$\varphi'(s)<0$ for small positive $s$.
Evenness also gives
\[
 \varphi(s)=\varphi(0)+\tfrac12\varphi''(0)s^2+O(s^4),\quad
 s^2=-\frac{2\delta}{\beta_c^2\varphi''(0)}+O(\delta^2),
\]
\[
 D=\varphi(0)s+\tfrac16\varphi''(0)s^3+O(s^5).
\]
Thus
\[
 \eta_* =\frac\delta{\beta_c}s+\frac{\beta_c}6\varphi''(0)s^3
 +O(\delta^{5/2})=\frac{2\delta}{3\beta_c}s+O(\delta^{5/2}),
 \qquad k_0=\frac\delta{3\beta_c}+O(\delta^2),
\]
proving the asymptotics. Directly,
$\eta_* =\beta\int_0^s[\varphi(q)-\varphi(s)]\,dq>0$.
For $|\eta|<\eta_*$, $0<k_\zeta^*<2k_0<1/2$.
The cubic coefficient decreases with $k$ and at its upper endpoint is
\[
 g_{mmm}|_{k=2k_0}=-\tfrac72\beta_c^3\varphi''(0)+o(1)>0.
\]
Proposition~\ref{p:sadCus} therefore makes both candidates
nondegenerate cusps unfolded by $(h,k)$.
At the boundary the couplings are $0,2k_0$: the first has $J=0$,
whereas the second retains its nonzero cubic coefficient and stable
complementary eigenvalue, which persist nearby.
\end{proof}

\paragraph{Remarks.}
The exact inequality $|\eta|<\beta(\Phi(s)-1/2)-s$, $\beta \Phi'(s)=1$,
concerns the two inherited cusps, not possible distant ones.
Their states and common stimulus remain fixed, while
$k_+^*-k_-^*=\eta/(\beta D)$. Since $D=O(\delta^{1/2})$ and
$k_0=O(\delta)$, stimulus asymmetry of order $\delta^{3/2}$ changes
the coupling by its own order. For the standard Gaussian response $\Phi_{\mathrm G}$,
\[
 \beta_c=\sqrt{2\pi},\qquad s=\sqrt{2\log(\beta/\beta_c)},
\]
\[
 \begin{aligned}
 \eta_* &=\beta\left[\Phi_{\mathrm G}\!\left(\sqrt{2\log(\beta/\beta_c)}\right)-\tfrac12\right]
 -\sqrt{2\log(\beta/\beta_c)}\\
 &=\frac{2\sqrt2}{3\beta_c^{3/2}}(\beta-\beta_c)^{3/2}
 +O((\beta-\beta_c)^{5/2}).
 \end{aligned}
\]
Unequal removal coefficients or asymmetric response shapes require
separate unfolding calculations; these formulas assume a symmetric
response and unequal additive stimuli.

\subsection{Odds-polynomial responses: a siphon-preserving cusp}
\label{ss:oddsCusp}

The constructions above give two distinct mechanisms: flat zero-onset
responses preserve siphons but have vanishing onset slope, whereas the
full-support responses in the explicit two-population cusp calculations
exclude nontrivial zero faces. Polynomial odds provide a common class in
which the response zero set is preserved while its onset slope and
nonlinear shape can be varied independently.

We use the cusp continuation criterion of Proposition~\ref{p:sadCus}
and persistence of hyperbolic equilibria under weak coupling
\cite{golubitsky2023dynamics,kuznetsov2004,Boud26}.
The constructions below additionally preserve specified siphon faces
and control their transverse stability.

\beD[Odds-polynomial threshold response]
[][Odd]
Fix $m\ge1$ and $a_1,\ldots,a_m\ge0$ with $a_m>0$, and put
\begin{equation}\label{eq:oddsFamily}
 \varphi(q)=\sum_{k=1}^m a_kq^k,\qquad
 \Phi(q)=\begin{cases}\dfrac{\varphi(q)}{1+\varphi(q)},&q\ge0,\\[4pt] 0,&q\le0.\end{cases}
\end{equation}
Then $\Phi$ is a nondecreasing distribution function, is strictly positive
for $q>0$, tends to one as $q\to\infty$, and has zero set $(-\infty,0]$.
It is smooth on $[0,\infty)$ from the right, with
\[
 \Phi'(0+)=a_1,\qquad \Phi''(0+)=2(a_2-a_1^2).
\]
If $a_1>0$, its extension by zero to negative arguments is not $C^1$ at
zero. Derivatives at zero below are taken in the admissible nonnegative
directions; local cusp calculations are performed at strictly
positive arguments, where the vector field is smooth. Since $1+\varphi>0$ on a
neighbourhood of $[0,\infty)$, the formula $\varphi/(1+\varphi)$ is a smooth
extension of $\Phi$ there, so threshold models built from $\Phi$ with nonnegative
polynomials $P_i$ satisfy the smoothness convention of Definition~\ref{d:KB}
(Remark~\ref{r:KBonset}).
\eeD

\beP[Odds responses keep the siphon lattice and scale the invasion matrix][odds invasion matrix][Odd]
Consider a production threshold model \eqref{eq:KB}, \eqref{polR} with $h_i=0$, a common
response $\Phi$, and polynomials with $P_i(0)=0$. Keep $V$, $b_i$, and $P_i$ fixed while varying the
admissible coefficients of $\Phi$ in \eqref{eq:oddsFamily}.
Write $z=x_\Sigma$ and $y=x_{\Sigma^c}$.
A nonempty set $\Sigma$ is a siphon if and only if
\[
 P_i(0,y)=0\quad(i\in\Sigma)
 \quad\text{for every admissible resident state }y,
\]
and
\begin{equation}\label{eq:oddsNoLinCoupling}
 V_{\Sigma\Sigma^c}=0.
\end{equation}
Consequently the entire siphon lattice is independent of the odds
coefficients. On such a face the transverse Jacobian is
\begin{equation}\label{eq:oddsTransverse}
 J_\Sigma^\perp(y)
 =-V_\Sigma
 +a_1\operatorname{Diag}(b_i:i\in\Sigma)D_zP_\Sigma(0,y).
\end{equation}
Thus only $a_1$ enters the transverse matrix at a \emph{fixed} resident
state $y$, and it changes that matrix only if
$D_zP_\Sigma(0,y)\ne0$. Higher odds coefficients may change invasion
thresholds indirectly, by moving the resident equilibrium at which
\eqref{eq:oddsTransverse} is evaluated.
\eeP
\begin{proof}
The polynomials $P_i$ are nonnegative on the state space, so there $\Phi$ may be replaced by
$\varphi/(1+\varphi)$, which is smooth near $q=0$; this gives a $C^1$ extension of the vector field,
with $\Pi\ge0$ on the admissible state space. Proposition~\ref{p:InvF}(1) then applies: $\Sigma$ is a siphon exactly when
$V_{\Sigma\Sigma^c}=0$ and $b_i\Phi(P_i(0,y))=0$ for $i\in\Sigma$. Since $b_i>0$ and $\Phi(q)>0$ for
$q>0$, the second condition is $P_i(0,y)=0$, which proves the equivalence and independence from
the response coefficients. Since $\Phi'(0+)=a_1$, the chain rule gives
$\partial_z\Pi_\Sigma(0,y)=a_1\operatorname{Diag}(b_i)D_zP_\Sigma(0,y)$, and the formula for
$J^\perp$ in Proposition~\ref{p:InvF}(1) proves \eqref{eq:oddsTransverse}.
\end{proof}

We next show why degree four, rather than a particular numerical choice
of coefficients, produces a positive cusp together with a stable
boundary equilibrium.

\beP[Quartic odds attain the minimum degree for a positive cusp above a stable origin][minimal degree and quartic odds cusp][Qua]\label{p:Min}
For $\dot x=-x+\Phi(x)$ with $\Phi$ defined by \eqref{eq:oddsFamily},
degree four is the minimum admitting both a hyperbolically stable
origin and a positive nondegenerate cusp. The following quartic
family attains it and also has a second hyperbolically stable equilibrium.

Choose $0<r<s<1$ with $3r+s\le1$, and define
\[
 c=\frac{1}{(1-r)^3(1-s)},\qquad
 \varphi_{r,s}(q)=\frac{q-cq(q-r)^3(q-s)}{1-q}.
\]
The apparent singularity at $q=1$ is removable: the expression is the
quartic polynomial
\[
 \begin{aligned}
 \varphi_{r,s}(q)&=a_1q+a_2q^2+a_3q^3+a_4q^4,\\
 a_1&=1-cr^3s,\\
 a_2&=c(1-3r-s+3r^2+3rs),\\
 a_3&=c(1-3r-s),\qquad a_4=c.
 \end{aligned}
\]
The coefficients $a_1,a_2,a_4$ are positive, and $a_3\ge0$, with equality
exactly when $3r+s=1$; thus $3r+s\le1$ is precisely the admissibility
condition $a_k\ge0$. With the response $\Phi_{r,s}$
constructed from $\varphi_{r,s}$, the scalar equation
$\dot x=f_{r,s}(x)=-x+\Phi_{r,s}(x)$ on $[0,1]$ satisfies
\[
 f_{r,s}(x)=-\frac{cx(x-r)^3(x-s)}{1+\varphi_{r,s}(x)}.
\]
It has a hyperbolically stable origin, a nondegenerate positive cusp
at $x=r$, and a hyperbolically stable positive equilibrium at $x=s$.
More precisely,
\[
 f_{r,s}'(0)=-cr^3s<0,\qquad
 f_{r,s}'''(r)=6cr(s-r)(1-r)>0,
\]
\[
 f_{r,s}'(s)=-cs(s-r)^3(1-s)<0.
\]
The cusp is generically unfolded by $(a_1,a_2)$, keeping $a_3,a_4$
fixed, since for any quartic odds polynomial
\begin{equation}\label{eq:oddsUnfold}
 \det\frac{\partial(f,f_x)}{\partial(a_1,a_2)}
 =\frac{x^2}{(1+\varphi(x))^4}>0\qquad(x>0).
\end{equation}
In particular, at $x=r$ the determinant equals $r^2(1-r)^4$.

The root parameters control the boundary stability margin and the cusp
strength,
\begin{equation}\label{eq:rsMargin}
 1-a_1=c\,r^3s,\qquad f_{r,s}'''(r)=6cr(s-r)(1-r),
\end{equation}
and over the admissible region
\begin{equation}\label{eq:rsSup}
 \sup_{0<r<s,\ 3r+s\le1}\,(1-a_1)=\frac1{81}.
\end{equation}
The supremum is not attained. It is approached only as $r,s\to1/4$,
where $f_{r,s}'''(r)\to0$ and the cusp degenerates. For instance, $(r,s)=(1/5,2/5)$ has $a_3=0$ and margin $1/96$.

For $r=1/8$ and $s=1/2$ this is exactly the reference polynomial
\[
 \varphi(q)=\tfrac1{343}(342q+368q^2+128q^3+1024q^4),
\]
for which
\begin{equation}\label{eq:oddsFactor}
 f(x)=-\frac{x(2x-1)(8x-1)^3}{343(1+\varphi(x))}.
\end{equation}
Here $f'(0)=-1/343$, $f'(1/2)=-27/686$, and
\[
 f=f_x=f_{xx}=0,\qquad f_{xxx}=\frac{36}{49}
 \quad\text{at }x=\frac18.
\]
At this cusp the unfolding determinant in \eqref{eq:oddsUnfold}
is $2401/262144$, and
\[
 \det\frac{\partial(f,f_x,f_{xx})}{\partial(x,a_1,a_2)}
 =\frac{441}{65536}\ne0
 \quad\text{at }(x,a_1,a_2)
 =\left(\tfrac18,\tfrac{342}{343},\tfrac{368}{343}\right).
\]
\eeP
\begin{proof}
\emph{Minimum degree.}
The numerator $(1-x)\varphi(x)-x$ has degree at most $m+1$,
and contains the factor $x$. A positive triple equilibrium thus
requires $m\ge3$. If $m=3$, its leading coefficient is $-a_3$
and necessarily
\[
 (1-x)\varphi(x)-x=-a_3x(x-r)^3,\qquad r>0.
\]
Comparison of the coefficients of $x$ yields
$a_1-1=a_3r^3>0$. Thus the origin is unstable.
For $m\le2$ the numerator cannot contain both $x$ and a
positive triple root. This proves the lower bound; the construction
below proves attainment.

\emph{Quartic construction.}
The choice of $c$ makes
$q-cq(q-r)^3(q-s)$ vanish at $q=1$, so division by $1-q$
gives a polynomial of degree four with the displayed coefficients.
Since $c>0$, $a_3\ge0$ is $3r+s\le1$, and
$a_2=a_3+3cr(r+s)>0$. Also $r<s\le1-3r$ gives $r<1/4$ and
$1-s\ge3r$, which imply $(1-r)^3(1-s)>r^3s$, hence $a_1>0$;
finally $a_4=c>0$.
The identity
$(1-x)\varphi_{r,s}(x)-x=-cx(x-r)^3(x-s)$ proves the
general factorization and hence its specialization
\eqref{eq:oddsFactor}. The three derivatives follow by factoring at
the respective equilibria and using
$1+\varphi_{r,s}(r)=(1-r)^{-1}$ and
$1+\varphi_{r,s}(s)=(1-s)^{-1}$.
For the unfolding determinant, differentiating
$f=(\varphi-x-x\varphi)/(1+\varphi)$ with respect to $a_1,a_2$
and then $x$ gives \eqref{eq:oddsUnfold}; at the cusp,
$f=f_x=f_{xx}=0$ and $f_{xxx}\ne0$. Hence the Jacobian
of $(f,f_x,f_{xx})$ in $(x,a_1,a_2)$ equals
$f_{xxx}$ times the unfolding determinant, and is nonzero. Proposition~\ref{p:sadCus} in one dimension,
therefore gives the claimed generically unfolded cusp.
Since $f_{r,s}'''(r)>0$, the cusp equilibrium itself is unstable;
the two stable equilibria are $0$ and $s$.
\emph{Stability margin.}
The margin identity in \eqref{eq:rsMargin} is $f_{r,s}'(0)=a_1-1$.
For \eqref{eq:rsSup}, the margin
$\Delta(r,s)=1-a_1=r^3s/((1-r)^3(1-s))$ increases with $s$, so
$\Delta(r,s)\le\Delta(r,1-3r)=r^2(1-3r)/(3(1-r)^3)$, and
\[
 \frac{d}{dr}\,\frac{r^2(1-3r)}{3(1-r)^3}=\frac{2r(1-4r)}{3(1-r)^4}>0
 \qquad(0<r<\tfrac14).
\]
The bound therefore increases to its value $1/81$ at $r=1/4$, which is not
admissible since there $s=1-3r=r$. Along any sequence with
$\Delta\to1/81$ both $r$ and $s$ tend to $1/4$, so $s-r\to0$ and
$f_{r,s}'''(r)\to0$.
Substitution gives the stated exact values.
\end{proof}

\paragraph{The quartic odds construction has an explicit scalar energy.}
For $\dot x=-x+\Phi_a(x)$ with an odds response $\Phi_a$ of Proposition~\ref{p:Qua}, $\Phi_a$ is strictly
increasing on $[0,\infty)$ when $a_1>0$, maps that interval onto $[0,1)$, and has $\Phi_a^{-1}(0)=0$. Define
\[
 \mathcal E_a(x)=\int_0^x\Phi_a^{-1}(s)\,ds-\frac{x^2}{2},\qquad 0\le x<1.
\]
Then $\dot{\mathcal E}_a(x)=-(\Phi_a^{-1}(x)-x)(x-\Phi_a(x))\le0$, with equality exactly at scalar equilibria.
At a positive equilibrium $e$, $\mathcal E_a''(e)=1/\Phi_a'(e)-1$, so a hyperbolically stable equilibrium
is a strict local minimum of this energy. At a scalar cusp $r$, where $\Phi_a(r)=r$, $\Phi_a'(r)=1$ and
$\Phi_a''(r)=0$, differentiating the inverse gives
\[
 \mathcal E_a'(r)=\mathcal E_a''(r)=\mathcal E_a'''(r)=0,\qquad \mathcal E_a''''(r)=-\Phi_a'''(r).
\]
For the reference quartic,
\[
 \mathcal E_a''(0+)=\frac1{342}>0,\qquad
 \mathcal E_a''(1/2)=\frac{27}{659}>0,\qquad
 \mathcal E_a''''(1/8)=-\frac{36}{49}<0.
\]
The stable origin and the stable positive equilibrium are therefore local energy minima, while the
cusp is a degenerate local maximum at the reference parameter. This is the generator of the Fenchel--Bregman divergence
of Subsection~\ref{s:FenBre} applied to the quartic construction, not an additional minimum-degree
claim. The coupled odds networks have nonlinear $P_i$, so this energy does not simply sum over them.

\beT[Networked odds responses carry an interior cusp over a stable siphon face]
[][Net]
Let $\Phi_a$ denote the odds response with $a_3=128/343$ and
$a_4=1024/343$ fixed, and initially
$(a_1,a_2)=a^0=(342/343,368/343)$ as in
Proposition~\ref{p:Qua}. Consider the two-population family on
the forward-invariant square $[0,1]^2$
\begin{equation}\label{eq:oddsNet}
 \dot x=-x+\Phi_a\bigl(x(1+\gamma y)\bigr),\qquad
 \dot y=-y+\Phi_a\bigl(y(10+\varepsilon x)\bigr),
 \qquad \gamma,\varepsilon\ge0.
\end{equation}
\BEN
\item For every admissible $a$ and every $\gamma,\varepsilon\ge0$,
 the faces $x=0$ and $y=0$ are invariant and the minimal siphons
 are exactly $\{x\}$ and $\{y\}$.
\item At $(\gamma,\varepsilon,a)=(0,0,a^0)$, the resident equation
 $\dot y=-y+\Phi_{a^0}(10y)$ has exactly one positive equilibrium
 $y^*=q^*/10$, where $q^*$ is the unique positive root of
\begin{equation}\label{eq:oddsResident}
 p(q)=1024q^4-10112q^3-912q^2-3338q-3077=0.
\end{equation}
 Its eigenvalue is strictly negative, and is given exactly by
 \[
 \lambda_y=-\frac{q^*p'(q^*)}{343(1+\varphi(q^*))}<0.
 \]
 In particular, at $(1/8,y^*)$ the full Jacobian has a simple
 zero eigenvalue and a strictly negative complementary eigenvalue.
\item There exists $\delta>0$ and smooth functions
 $a_1(\gamma,\varepsilon)$, $a_2(\gamma,\varepsilon)$,
 and $x_c(\gamma,\varepsilon)$, for
 $0\le\gamma,\varepsilon<\delta$, reducing to
 $(a_1^0,a_2^0,1/8)$ at zero coupling, such that
 \eqref{eq:oddsNet} has a nondegenerate interior cusp near
 $(1/8,y^*)$. Its complementary eigenvalue remains strictly
 negative, and $(a_1,a_2)$ still unfold the cusp generically.
\item Let $y_*(a_1,a_2)$ be the unique resident equilibrium
 near $y^*$ for the current coefficients. On the continued cusp
 family, the boundary equilibrium $(0,y_*)$ has transverse eigenvalue
\begin{equation}\label{eq:oddsTransEig}
 \lambda_{\mathrm{tra}}(\gamma,\varepsilon)
 =-1+a_1(\gamma,\varepsilon)
 \Bigl(1+\gamma y_*\bigl(a_1(\gamma,\varepsilon),
                             a_2(\gamma,\varepsilon)\bigr)\Bigr).
\end{equation}
 This is $-1/343$ at zero coupling and remains negative for all
 sufficiently small nonnegative couplings.
\item Along the continued cusp family, a stable boundary
 equilibrium near $(0,y^*)$ and a stable interior equilibrium
 near $(1/2,y^*)$ coexist. For $\gamma,\varepsilon>0$, both off-diagonal Jacobian entries
 are strictly positive at every positive interior state.
\EEN
Consequently, sufficiently small positive coupling preserves
both minimal siphons, a nondegenerate interior cusp, and bistability
between a stable boundary equilibrium and a stable positive equilibrium,
with a strictly stable transverse invasion eigenvalue on the boundary.
\eeT
\begin{proof}
For part 1, both arguments of $\Phi_a$ contain the corresponding
variable as a factor. The zero-onset property therefore makes each
coordinate zero face invariant. These singleton faces are the
minimal siphon faces.

For part 2, with $q=10y$ one computes
\[
 -y+\Phi_{a^0}(10y)
 =-\frac{q\,p(q)}{10\,[343+343\varphi(q)]}.
\]
The coefficients of $p$ have exactly one sign change, so by
Descartes' rule $p$ has exactly one positive root, counted with
multiplicity. Since $p(0)<0$ and $p(q)\to\infty$, that root is
simple and $p'(q^*)>0$. Differentiating the displayed identity at
$q=q^*$ gives $\lambda_y$ as stated. The equations decouple at
zero coupling, and Proposition~\ref{p:Qua} gives the zero eigenvalue
in the $x$ direction.

For part 3, at zero coupling the first equation is the scalar
cusp of Proposition~\ref{p:Qua}, unfolded by $(a_1,a_2)$; the
second equation is independent of $x$ and has eigenvalue
$\lambda_y<0$. Eliminating its equilibrium equation therefore
leaves exactly $\Phi_a(x)-x$, so the full system satisfies the cusp
tests of Proposition~\ref{p:sadCus}. Its persistence assertion
continues the cusp for small $(\gamma,\varepsilon)$, preserving
the nonzero cubic, the unfolding determinant, and the stable
complementary eigenvalue. Smoothness of this family gives the
claimed smooth continuation. The base state and odds coefficients
are strictly positive; restricting the continued family to
nonnegative couplings preserves admissibility.

For part 4, restriction to $x=0$ gives
$\dot y=-y+\Phi_a(10y)$, independent of $\gamma,\varepsilon$.
Its hyperbolic positive equilibrium continues smoothly with $a$.
The face Jacobian is triangular, with transverse entry
$-1+\Phi_a'(0+)(1+\gamma y_*)$, which is
\eqref{eq:oddsTransEig}. This is strictly negative at the base
point, so it remains so for sufficiently small couplings.

For part 5, both $(0,y^*)$ and $(1/2,y^*)$ are hyperbolically
stable at zero coupling: their tangential resident eigenvalue is
$\lambda_y<0$, the former has transverse eigenvalue $-1/343$,
and the latter has $x$ eigenvalue $-27/686$. The implicit-function
theorem and spectral continuity preserve both equilibria and their
stability under the continued parameter adjustment. The first stays
on the invariant face $x=0$; the second remains strictly positive.
For positive couplings, direct differentiation gives
$\partial_y\dot x=\gamma x\Phi_a'(x(1+\gamma y))>0$
and
$\partial_x\dot y=\varepsilon y\Phi_a'(y(10+\varepsilon x))>0$
when $x,y>0$.
\end{proof}

\beR[What the construction does and does not settle]
For fixed odds coefficients from Proposition~\ref{p:Qua}, let $y^*$
be a positive equilibrium of the corresponding resident equation
$\dot y=-y+\Phi_{r,s}(10y)$. The transverse stability condition at
$(0,y^*)$ is, by \eqref{eq:rsMargin},
\begin{equation}\label{eq:rsGamma}
 \gamma<\frac{1-a_1}{a_1y^*}=\frac{c\,r^3s}{(1-c\,r^3s)\,y^*},
\end{equation}
which is $-1+a_1(1+\gamma y^*)<0$ solved for $\gamma$. For the reference
coefficients it is $\gamma<1/(342y^*)\approx2.9\times10^{-3}$.
This is \emph{not} an explicit bound along the continued cusp family: there $a_1$ changes with
the couplings, and within the quartic family changing $(r,s)$ also changes the resident state,
the cusp location and its cubic coefficient. The supremum $1/81$ in \eqref{eq:rsSup} therefore
bounds the uncoupled origin's stability margin, not the admissible coupling. The two-parameter quartic
family of Proposition~\ref{p:Qua} permits the cusp location, the
second stable equilibrium, and the boundary stability margin to be
varied while the response zero set is fixed; its coupled continuation
requires only the same hyperbolicity and unfolding hypotheses.
The theorem gives local continuation, not a classification of all cusp
branches by their invasion thresholds.
\eeR

\beR[Why the coupling must be multiplicative for singleton siphons][multiplicative coupling][Mul]
The arguments of \eqref{eq:oddsNet} are $x(1+\gamma y)$ and $y(10+\varepsilon x)$,
not $x+\gamma y$ and $10y+\varepsilon x$. For a zero-onset response that is strictly positive on $(0,\infty)$,
diagonal removal and nonnegative polynomials $P_i$, the singleton $\{i\}$ is a
siphon exactly when
\[
 P_i(x)\big|_{x_i=0}\equiv0,
\]
that is, when every monomial of $P_i$ contains $x_i$; this is the criterion of
Proposition~\ref{p:KBsiph} specialised to a singleton. The additive argument
$x+\gamma y$ has the monomial $\gamma y$, which survives at $x=0$, so
$\Phi(\gamma y)>0$ and the face $x=0$ is not invariant: the siphon is destroyed for
every $\gamma>0$. The multiplicative argument $x(1+\gamma y)$ carries $x$ in every
monomial and so preserves it for all $\gamma$. A nonzero off-diagonal linear term therefore destroys the
corresponding singleton siphon.

The necessity is specific to singleton siphons. For a general $\Sigma$, with
$V_{\Sigma\Sigma^c}=0$, the criterion of Proposition~\ref{p:Odd} requires only
\[
 \operatorname{supp}\nu\cap\Sigma\ne\varnothing
 \quad\text{for every monomial }x^\nu\text{ of }P_i,\ i\in\Sigma,
\]
so a coordinate of $\Sigma$ may have an additive term from another coordinate of
$\Sigma$. Overlapping minimal siphons therefore allow a richer class of couplings.
\eeR

% Former separate minimum-degree proposition (merged above; inactive).
% \beP[Quartic is the minimal degree carrying a stable origin and a positive cusp][minimal odds degree][Min]
% For the scalar equation $\dot x=f(x)=-x+F(x)$ with $F$ as in
% \eqref{eq:oddsFamily} and $a_k\ge0$, no degree $m\le3$ admits
% simultaneously a hyperbolically stable origin and a positive
% nondegenerate cusp. The quartic family of
% Proposition~\ref{p:Qua} therefore attains the minimum degree.
% \eeP
% \begin{proof}
% The numerator $(1-x)\phi(x)-x$ has degree at most $m+1$,
% and contains the factor $x$. A positive triple equilibrium thus
% requires $m\ge3$. If $m=3$, its leading coefficient is $-a_3$
% and necessarily
% \[
%  (1-x)\phi(x)-x=-a_3x(x-r)^3,\qquad r>0.
% \]
% Comparison of the coefficients of $x$ yields
% $a_1-1=a_3r^3>0$. Thus the origin is unstable.
% For $m\le2$ the numerator cannot contain both $x$ and a
% positive triple root.
% \end{proof}
%

% Former longer margin discussion (inactive).
% \beR[The margin trade-off is not monotone in $a_1$][margin trade-off][Tra]
% Lowering $a_1$ does improve
% \eqref{eq:rsGamma} with everything else fixed, but $a_1$ is not a free
% parameter: it is determined by $(r,s)$ through \eqref{eq:rsMargin}, and moving
% $(r,s)$ also moves the cusp location $r$, the second equilibrium $s$, and the
% cusp strength $f'''(r)$. The margin is therefore traded against the geometry of
% the equilibria, and the trade is constrained by $a_3\ge0$. By
% \eqref{eq:rsSup} the margin never reaches $1/81$ and approaches it only as
% the cusp degenerates. A quantitative coupling window along the continued cusp
% family is a separate problem, which remains open.
% \eeR 